\anexo{Brainstorming: material de la sesión}
\label{anx:brainstorming}


Este anexo presenta el material fotográfico y la tabla completa de ideas producidas durante la sesión de brainstorming del 10 de octubre de 2026, descripta en la sección~\ref{sec:brainstorming}.

\section*{\thechapter.1 Fotografías de la sesión}
\addcontentsline{toc}{subsection}{\thechapter.1 Fotografías de la sesión}

\begin{figure}[ht]
    \centering
    \includegraphics[width=0.85\textwidth]{./images/brainstoming-2.jpeg}
    \caption{Agrupación por afinidad y votación por puntos: ideas organizadas por impacto percibido (sesión del 10 de octubre de 2026).}
    \label{fig:brainstorming-grupos}
\end{figure}

\noindent\textit{Fuente: elaboración propia.}

\FloatBarrier

\begin{figure}[ht]
    \centering
    \includegraphics[width=\textwidth]{./images/brainstoming-3.jpeg}
    \caption{Matriz impacto/esfuerzo con los cuatro cuadrantes: ideas más votadas ubicadas según impacto percibido (eje vertical) y esfuerzo técnico estimado (eje horizontal). Sesión del 10 de octubre de 2026.}
    \label{fig:brainstorming-matriz-completa}
\end{figure}

\noindent\textit{Fuente: elaboración propia.}

\FloatBarrier

\begin{figure}[ht]
    \centering
    \includegraphics[width=0.75\textwidth]{./images/brainstoming-1.jpeg}
    \caption{Detalle del sector de alto esfuerzo de la matriz impacto/esfuerzo (sesión del 10 de octubre de 2026).}
    \label{fig:brainstorming-matriz-detalle}
\end{figure}

\noindent\textit{Fuente: elaboración propia.}

\FloatBarrier

\section*{\thechapter.2 Tabla completa de ideas}
\addcontentsline{toc}{subsection}{\thechapter.2 Tabla completa de ideas}

La Tabla~\ref{tab:brainstorming-ideas} presenta las 21 ideas producidas en la sesión, con su cuadrante en la matriz impacto/esfuerzo y su estado de implementación al cierre del proyecto.

\begin{table}[ht]
\centering
\caption{Ideas de la sesión de brainstorming y estado de implementación.}
\label{tab:brainstorming-ideas}
\begin{tabular}{p{5.5cm} p{3cm} p{5.5cm}}
\hline
\textbf{Idea} & \textbf{Cuadrante} & \textbf{Estado} \\
\hline
\multicolumn{3}{l}{\textit{Ideas que validan el diseño implementado}} \\
\hline
Pedir feedback al usuario & Alto / poco esfuerzo & Implementada: marcar como ejecutada, deshacer, «No fui yo» (HU-07). \\
Efemérides: avisos de fechas clave (Black Friday, Día de la Madre) & Alto / poco esfuerzo & Implementada: recomendación de fecha clave con calendario por rubro (HU-06). \\
Efemérides: qué productos promocionar & Alto / medio esfuerzo & Implementada: fecha clave sobre los productos más vendidos de los últimos 28 días. \\
Promos & Alto / poco esfuerzo & Implementada: rotación en caída y fechas clave sugieren promociones. \\
Segmentar clientes en alto, medio y bajo & Alto / poco esfuerzo & Implementada: segmentación RFM con \emph{k}-means y CLV (HU-10). \\
Explicar por qué la recomendación & Alto / medio esfuerzo & Implementada: campo «Por qué te lo recomendamos» y explicación de la priorización en el detalle. \\
Productos estacionales & Medio / mucho esfuerzo & Implementada: pronóstico Holt-Winters con estacionalidad y preparación anticipada (HU-02). \\
Mostrar casos de éxito & Alto / poco esfuerzo & Implementada: vista de impacto esperado versus real (HU-08). \\
Detectar si el emprendedor implementó la recomendación & Alto / poco esfuerzo & Implementada: inferencia automática de la ejecución a partir de snapshots de stock y precio (HU-07). \\
\hline
\multicolumn{3}{l}{\textit{Ideas parcialmente implementadas}} \\
\hline
Diferenciar notificaciones de recomendaciones & Alto / poco esfuerzo & Parcial: niveles urgente, atención y oportunidad, sin separación visual entre avisos y sugerencias. \\
Pop-up para preguntar si le sirvió & Alto / poco esfuerzo & Parcial: impacto medido automáticamente; falta consulta activa al emprendedor. \\
Botón para realizar la acción desde la tarjeta & Alto / medio esfuerzo & Parcial: CTA en cada tarjeta; sin integración de escritura hacia la plataforma (trabajo futuro: «acciones con un clic»). \\
Dar solo recomendaciones que generen buen margen & Alto / mucho esfuerzo & Parcial: el bandit prioriza lo que funcionó; falta filtro por margen (requiere el costo del producto). \\
No hablar en imperativo & Alto / poco esfuerzo & Parcial: el LLM redacta los mensajes; aplicación efectiva depende del \emph{system prompt} de Gemini. \\
Adaptar el tono al rubro & Alto / poco esfuerzo & Parcial: el LLM recibe el rubro como contexto; aplicación efectiva depende del \emph{system prompt} de Gemini. \\
\hline
\multicolumn{3}{l}{\textit{Trabajo futuro}} \\
\hline
Mostrar un estado general (pérdidas y ganancias acumuladas) & Alto / mucho esfuerzo & Trabajo futuro: viable solo como impacto acumulado de recomendaciones; sin ese encuadre replica dashboards descriptivos. \\
Feedback de acciones ejecutadas fuera de la plataforma & Alto / mucho esfuerzo & Trabajo futuro: no observable desde la API de la plataforma. \\
Productos en tendencia del rubro que la tienda no tiene & Alto / mucho esfuerzo & Trabajo futuro: requiere fuentes de datos de mercado externas, fuera del alcance actual. \\
Enviar reporte mensual & Bajo / poco esfuerzo & Trabajo futuro: bajo impacto percibido en la matriz. \\
Dar recomendaciones al cierre de cada semana & Bajo / poco esfuerzo & Trabajo futuro: bajo impacto percibido en la matriz. \\
\hline
\multicolumn{3}{l}{\textit{Descartadas}} \\
\hline
Cambio de tendencia en reseñas negativas & Bajo / mucho esfuerzo & Descartada: la API de Tiendanube no expone reseñas de forma nativa; impacto percibido bajo. \\
\hline
\end{tabular}
\end{table}

\noindent\textit{Fuente: elaboración propia a partir de los resultados de la sesión del 10 de octubre de 2026.}
